\documentclass[11pt,a4paper]{article}
\usepackage[margin=0.55in, top=0.5in, bottom=0.5in]{geometry}
\usepackage[utf8]{inputenc}
\usepackage[T1]{fontenc}
\usepackage{charter}
\usepackage{fontawesome5}
\usepackage{xcolor}
\usepackage{titlesec}
\usepackage{enumitem}
\usepackage{tikz}
\usepackage{hyperref}
\usepackage{tcolorbox}
\tcbuselibrary{skins, breakable}
\usepackage{microtype}

\definecolor{headerbg}{HTML}{FFFFFF}
\definecolor{accent}{HTML}{1E3A5F}     % Professional Dark Navy
\definecolor{accentblue}{HTML}{2B6CB0} % Slate Blue
\definecolor{lightgray}{HTML}{F3F4F6}
\definecolor{darktext}{HTML}{1F2937}
\definecolor{skillbg}{HTML}{374151}
\definecolor{bulletline}{HTML}{9CA3AF}
\definecolor{subtletext}{HTML}{4B5563}

\hypersetup{colorlinks=true, linkcolor=accentblue, urlcolor=accentblue}

\titleformat{\section}{\large\bfseries\color{accent}}{}{0em}{}[\vspace{1pt}\titlerule]
\titlespacing{\section}{0pt}{11pt}{7pt}

\newtcbox{\skillbadge}{
    on line, box align=base,
    colback=skillbg, colframe=skillbg,
    colupper=white,
    fontupper=\small\bfseries,
    arc=4pt, boxrule=0pt,
    left=6pt, right=6pt, top=3pt, bottom=3pt,
    boxsep=0pt
}

\newtcolorbox{timelineentry}{
    blanker,
    breakable,
    borderline west={1.5pt}{0pt}{bulletline},
    left=1.2em,
    top=0pt,
    bottom=0pt,
    before skip=0.8ex,
    after skip=2.6ex
}

\newcommand{\cventry}[5]{%
  \noindent
  \begin{minipage}{\textwidth}
    \vspace{0pt}%
    \tikz[baseline=-0.7ex]{\node[draw=accent, fill=white, circle, inner sep=2.5pt, line width=1.5pt] {};}%
    \hspace{0.5em}%
    {\large\bfseries\color{accent} #1}%
    \hfill{\small\bfseries\color{accentblue} #4}\par
    \vspace{0.2ex}%
    \hspace{1.5em}{\normalsize\bfseries\color{darktext} #2}%
    \hfill{\small\color{subtletext}\textit{#3}}\par
  \end{minipage}\par
  \vspace{0.4ex}%
  \begin{timelineentry}
    \small\color{darktext}
    \begin{itemize}[label={\color{accentblue}\raisebox{0.2ex}{\tiny\faCircle}}, leftmargin=1.2em, itemsep=0.4ex, parsep=0pt, topsep=0pt, partopsep=0pt]
      #5
    \end{itemize}
  \end{timelineentry}
}

\begin{document}
\thispagestyle{empty}
\color{darktext}

% Encabezado
\begin{center}
    {\Huge \bfseries\color{accent} SANTIAGO ALEJANDRO OCHOA QUESADA} \\[6pt]
    {\small\color{darktext}
    \faEnvelope\ \href{mailto:santy.alejo4@gmail.com}{santy.alejo4@gmail.com}
    \quad\textcolor{bulletline}{\textbar}\quad
    \faPhone\ +57 320 351 7153
    \quad\textcolor{bulletline}{\textbar}\quad
    \faMapMarker*\ Bogotá, Colombia \\[4pt]
    \faLinkedin\ \href{https://linkedin.com/in/santiago-alejandro-ochoa-quesada-a07a18288/}{LinkedIn}
    \quad\textcolor{bulletline}{\textbar}\quad
    \faGithub\ \href{https://github.com/Sochoaq}{GitHub}
    }
\end{center}
\vspace{4pt}

\section{Perfil Profesional}
\textbf{\textbf{\textbf{\textbf{Ingeniero Electrónico}}}} de la \textbf{\textbf{\textbf{\textbf{Universidad Nacional de Colombia}}}} (promedio 4.6/5.0) y actual \textbf{\textbf{\textbf{\textbf{Ingeniero de Proyectos}}}} en AM4. Especializado en la ejecución técnica de proyectos de ingeniería industrial, montajes, pruebas de campo y supervisión de equipos de proceso y maquinaria crítica. Experiencia directa en instrumentación, pruebas y diagnóstico a motores y variación de frecuencia para clientes industriales.

\section{Experiencia Profesional}

\cventry{Ingeniero de Proyectos (Control y Automatización)}{AM4}{Bogotá}{04/2026 - Actualidad}
{
    \item Diseñé e implementé automatizaciones en Python y C/C++ para captura y procesamiento de variables industriales, facilitando el monitoreo continuo.
    \item Supervisé la resolución estructurada de problemas e inconsistencias en campo, coordinando proveedores y contratistas para el cumplimiento de cronogramas y especificaciones técnicas.
    \item Desarrollo de flujos para adquisición, limpieza y transformación de datos operativos provenientes de sistemas de medición, integrando diferentes fuentes de información.
    \item Implementación de automatizaciones en Python y C/C++ para capturar, organizar y procesar variables de operación industrial, reduciendo tareas manuales y facilitando el monitoreo.
    \item Ejecución de pruebas, validación y diagnóstico a motores y equipos industriales en proyectos para clientes como Refoenergy y Brinsa, asegurando parámetros de operación y confiabilidad.
    \item Realización de medidas técnicas e instrumentación en campo en proyectos para Geopark (DPS), garantizando la precisión, consistencia y calidad de los datos operativos.
    \item Programación básica de controladores lógicos programables (PLC) mediante lenguajes de descripción de hardware (HDL), apoyando la automatización de procesos industriales.
}

\cventry{Auxiliar Estudiantil}{Universidad Nacional de Colombia}{Bogotá}{07/2024 - 11/2025}
{
    \item Apoyo en laboratorios mediante organización, validación y análisis de resultados experimentales, programación, simulación y depuración de sistemas digitales y señales.
    \item Preparación y verificación de prácticas, registros, equipos de instrumentación y material técnico, asegurando consistencia en la toma de datos durante las sesiones.
    \item Acompañamiento a estudiantes en análisis de resultados, identificación de errores y resolución de problemas de electrónica y control durante prácticas de laboratorio.
    \item Elaboración de guías, reportes y presentaciones técnicas para documentar procedimientos, fallas, hallazgos y aprendizajes de manera estructurada.
}

\section{Educación}

\cventry{Ingeniería Electrónica}{Universidad Nacional de Colombia}{Bogotá}{02/2021 - 11/2025}
{
    \item \textbf{Promedio académico: 4.6 / 5.0}.
    \item Formación en control continuo y discreto, automatización, sistemas digitales, procesamiento de señales, instrumentación, comunicaciones y análisis numérico.
    \item Desarrollo de proyectos enfocados en simulación de procesos, programación de hardware, validación experimental y elaboración de reportes técnicos de ingeniería.
}

\newpage

\section{Proyectos Relevantes}

\cventry{Ingeniero de Proyectos (Control y Automatización)}{Ingeniería de Proyectos y Analítica}{Bogotá}{2026}
{
    \item Suite de instrumentación analítica para diagnóstico predictivo en campo bajo normativas IEEE e IEC, garantizando comisionamiento y alta disponibilidad operativa.
    \item Implementación de flujos de análisis, validación de señales y automatización de reportes operativos mediante Python y SQL.
    \item Desarrollo de componentes para almacenamiento, consulta, visualización y documentación de los resultados, facilitando la trazabilidad operativa.
}

\cventry{Procesamiento y análisis de sistemas de medición}{Ingeniería de Proyectos y Analítica}{Bogotá}{2026}
{
    \item Desarrollo de flujo automatizado en Python para adquirir, calibrar y validar datos de sensores industriales en tiempo real.
    \item Implementación de flujos de análisis, validación de señales y automatización de reportes operativos mediante Python y SQL.
    \item Desarrollo de componentes para almacenamiento, consulta, visualización y documentación de los resultados, facilitando la trazabilidad operativa.
}

\section{Habilidades}
\begin{tcolorbox}[colback=lightgray, colframe=lightgray, boxrule=0pt, arc=6pt, left=10pt, right=10pt, top=10pt, bottom=10pt]
\noindent
\linespread{1.8}\selectfont
\skillbadge{Automatización de Procesos}
\skillbadge{Control Industrial}
\skillbadge{Ciberseguridad y Redes}
\skillbadge{Cisco CCNA}
\skillbadge{Ethical Hacking}
\skillbadge{Linux \& Infraestructura}
\skillbadge{Python}
\skillbadge{C/C++}
\skillbadge{SQL}
\skillbadge{Adquisición de Datos}
\skillbadge{Validación de Señales}
\skillbadge{Gestión de Proyectos}
\skillbadge{Documentación Técnica}
\skillbadge{Resolución de Problemas}
\skillbadge{Trabajo en Equipo}
\end{tcolorbox}

\section{Idiomas}
\begin{itemize}[label={}, leftmargin=0.3cm, itemsep=6pt]
    \item \textbf{Español} \hfill Nativo \quad \textcolor{accent}{\faCircle \faCircle \faCircle \faCircle \faCircle}
    \item \textbf{Inglés} \hfill B2 (Profesional de Trabajo) \quad \textcolor{accent}{\faCircle \faCircle \faCircle \faCircle} \textcolor{bulletline}{\faCircle}
    \item \textbf{Francés} \hfill A2 \quad \textcolor{accent}{\faCircle \faCircle} \textcolor{bulletline}{\faCircle \faCircle \faCircle}
\end{itemize}

\section{Certificaciones}
\begin{itemize}[label={\color{accentblue}\faAward}, leftmargin=0.6cm, itemsep=6pt]
    \item \textbf{Ethical Hacking \& Cybersecurity} --- Cisco Networking Academy \hfill \textit{04/2026}
    \item \textbf{CCNA: Introduction to Networks} --- Cisco Networking Academy \hfill \textit{07/2025}
    \item \textbf{Análisis Computacional de Datos Ambientales} --- Universidad Nacional de Colombia \hfill \textit{09/2025}
    \item \textbf{Introducción a la Fotónica Integrada y las Comunicaciones Cuánticas} --- UNAL \hfill \textit{07/2024}
    \item \textbf{Python Intermedio} --- SoloLearn \hfill \textit{03/2021}
\end{itemize}

\end{document}
